\section{有限自动机识别正则语言}
\label{sec:04-Discrete-Mathematics/07-Boolean-Algebra-and-Automata/02}

有人提交了一个正则表达式，用来检查配置文件里的括号是否配平。评审意见只有一句：正则做不到。提交者当成是写法不够巧，换了三版，每一版都被一个更深的嵌套打穿。

这句``做不到''不是经验之谈。它有证明，而且证明用到的全部素材只有一句话：机器的状态数是有限的。本节要把这句话变成一件能反复使用的判断工具——先看有限状态能做什么，再看它做不到什么，以及边界画在哪里。

先看能做的那一侧。判断一个二进制串是否以 \(01\) 结尾，笨办法是把整串存下来，读完再回头看最后两位。可即使串长一百万位，影响答案的只有两件事：末位是不是 \(0\)（有可能充当 \(01\) 的开头），以及最后两位是不是刚好凑成 \(01\)。其余信息读过就可以扔。

\subsection{状态是对已读前缀的一句断言}

确定有限自动机（DFA）把这种``只留必要摘要''的做法固定下来。它由有限状态集 \(Q\)、输入字母表 \(\Sigma\)、转移函数 \(\delta:Q\times\Sigma\to Q\)、初始状态 \(q_0\) 和接受状态集 \(F\subseteq Q\) 组成。读入 \(w=a_1\cdots a_n\) 时从 \(q_0\) 出发反复代入 \(q_k=\delta(q_{k-1},a_k)\)，末状态落在 \(F\) 里就接受，否则拒绝。机器没有别的存储；\(\delta\) 的定义域是有限集，能记住的东西被 \(\abs{Q}\) 一次封顶。

语言在这里就是字符串集合 \(L\subseteq\Sigma^*\)，\(\Sigma^*\) 表示全体有限串（含空串 \(\varepsilon\)）。若某个 DFA 恰好接受 \(L\)，就称 \(L\) 正则。这个词有精确所指——能被有限记忆识别——与日常说的``看起来有规律''不是一回事。

\begin{center}
\begin{tikzpicture}[x=1cm,y=1cm,font=\small,>=stealth]
  \draw[fill=blue!6] (0,0) circle (0.5);
  \node at (0,0) {\(q_0\)};
  \draw[fill=blue!6] (3,0) circle (0.5);
  \node at (3,0) {\(q_1\)};
  \draw[fill=green!14] (6,0) circle (0.5);
  \draw (6,0) circle (0.62);
  \node at (6,0) {\(q_2\)};
  \draw[->,thick] (-1.4,0) -- (-0.55,0);
  \node[left,font=\footnotesize] at (-1.45,0) {开始};
  \draw[->] (0.52,0) -- (2.46,0);
  \node[above,font=\footnotesize] at (1.5,0.06) {\(0\)};
  \draw[->] (3.52,0) -- (5.46,0);
  \node[above,font=\footnotesize] at (4.5,0.06) {\(1\)};
  \draw[->] (0.20,0.474) arc (-60:240:0.40);
  \node[font=\footnotesize] at (0,1.45) {\(1\)};
  \draw[->] (3.20,0.474) arc (-60:240:0.40);
  \node[font=\footnotesize] at (3,1.45) {\(0\)};
  \draw[->] (5.59,-0.29) to[bend left=30] (3.41,-0.29);
  \node[font=\footnotesize] at (4.5,-0.86) {\(0\)};
  \draw[->] (6,-0.52) -- (6,-1.55) -- (0,-1.55) -- (0,-0.52);
  \node[below,font=\footnotesize] at (3,-1.58) {\(1\)};
\end{tikzpicture}

\emph{识别``以 \(01\) 结尾''的三状态机。\(q_0\) 记``末尾对 \(01\) 毫无进展''，\(q_1\) 记``末位是 \(0\)''，双圈的 \(q_2\) 记``最后两位正是 \(01\)''。任意长度的输入都被折进这三格之一，每读一个符号只更新一次。\(q_2\) 上读到 \(1\) 会退回 \(q_0\)：接受状态不是终点站，只是``此刻停下来的话答案为是''。}
\end{center}

拿 \(11001\) 走一遍：状态依次是 \(q_0\to q_0\to q_0\to q_1\to q_1\to q_2\)，接受。真正要带走的话写在图注里：把机器带到同一状态的两个前缀，此后的命运完全一致。状态不是流程图上随手画的圆圈，它是一句断言——``这一批前缀，我今后不再区分。''

\subsection{哪些前缀可以不再区分}

这句断言可以脱离机器单独定义。对语言 \(L\) 和两个前缀 \(x,y\)，若对一切后缀 \(z\) 都有

\[
xz\in L\iff yz\in L,
\]

就说 \(x\) 与 \(y\) 不可区分。反过来，只要存在一个后缀 \(z\) 让一边接受、另一边拒绝，任何识别 \(L\) 的机器都必须给它们安排不同状态——否则读完 \(z\) 之后两条输入会得到同一个答案。

\begin{center}
\begin{tikzpicture}[x=1cm,y=1cm,font=\small,>=stealth]
  \node[right] at (0,2.4) {\(\varepsilon,\ 1,\ 11,\ 0111,\ \dots\)};
  \node[right] at (0,1.2) {\(0,\ 10,\ 100,\ 0110,\ \dots\)};
  \node[right] at (0,0) {\(01,\ 101,\ 11001,\ \dots\)};
  \draw[->] (4.2,2.4) -- (5.3,2.4);
  \draw[->] (4.2,1.2) -- (5.3,1.2);
  \draw[->] (4.2,0) -- (5.3,0);
  \draw[fill=blue!8,rounded corners=3pt] (5.4,2.05) rectangle (10.4,2.75);
  \node at (7.9,2.4) {\(q_0\)：对 \(01\) 无进展};
  \draw[fill=blue!8,rounded corners=3pt] (5.4,0.85) rectangle (10.4,1.55);
  \node at (7.9,1.2) {\(q_1\)：末位是 \(0\)};
  \draw[fill=green!14,rounded corners=3pt] (5.4,-0.35) rectangle (10.4,0.35);
  \node at (7.9,0) {\(q_2\)：末两位是 \(01\)};
  \node[font=\footnotesize] at (2.1,-1.0) {无限多前缀};
  \node[font=\footnotesize] at (7.9,-1.0) {三个等价类};
\end{tikzpicture}

\emph{左边这一列永远列不完，右边只有三格。归并的依据不是长相相似，而是未来相同：同格的两个前缀，接上任何后缀之后接受与否始终一致。有限自动机之所以可能，就因为这种归并留下的类别只有有限个。}
\end{center}

这层关系叫 Myhill--Nerode 关系，它给出一句可以直接当判据用的话：等价类的个数有限，当且仅当 \(L\) 正则；而且最小 DFA 的状态数恰好等于类的个数。它比``先画一台机器再做最小化''更靠前，因为整句话里没有出现任何一台具体机器——类别是语言自己的性质，机器只是把类别实现出来。

用它算下界很方便。对``以 \(01\) 结尾''，取 \(\varepsilon\)、\(0\)、\(01\) 三个前缀：接上后缀 \(1\)，只有 \(0\) 那一支被接受；接上空串，只有 \(01\) 那一支被接受。三者两两可区分，所以任何识别它的 DFA 至少要三个状态，前面那台已经最小。

\subsection{非确定性压缩描述，不扩张能力}

非确定有限自动机（NFA）把转移放宽成 \(\delta:Q\times(\Sigma\cup\set{\varepsilon})\to\mathcal P(Q)\)：一次输入可以同时进入多个状态，还能沿 \(\varepsilon\)-边不消耗符号地跳。只要存在一条计算路径在读完输入时落在接受态，就算接受。它看起来像是允许机器``猜''——识别含子串 \(01\) 的语言时，一边留在扫描状态，一边在读到 \(0\) 时分出一条等待 \(1\) 的支线。

能力却没变。子集构造把 NFA 的一组可能状态 \(S\subseteq Q\) 当作 DFA 的单个状态，转移取 \(\bigcup_{q\in S}\delta(q,a)\) 再补上 \(\varepsilon\)-闭包，运行时维护的正是这个集合。状态数则最坏要涨到 \(2^{\abs{Q}}\)。非确定性买到的是描述长度，不是识别范围。

正则表达式给同一类语言换了一种写法：从空语言、空串和单符号出发，用并、连接与 Kleene 星 \(L^*=\bigcup_{k\ge0}L^k\) 往上搭，例如 \((0\lor1)^*01\) 就是本节那个语言。Kleene 定理说这两种描述覆盖的语言完全相同，表达式适合声明模式，自动机适合逐符号执行。至于闭包性质，直接由机器的组合给出：让两台 DFA 并行跑在乘积状态 \(Q_1\times Q_2\) 上，要求两个分量都接受就得到交，要求至少一个接受就得到并，把一台完整 DFA 的接受与拒绝状态对调就得到补。构造顺带给出了算法，状态数则要相乘。

\subsection{有限内存做不到的事}

现在换到另一侧。考虑

\[
L=\set{a^nb^n:n\ge0},
\]

也就是若干个 \(a\) 后面跟着同样多的 \(b\)。要判断两边数量相等，机器读完 \(a\) 那一段时必须知道读了多少个；\(n\) 没有上界，而状态只有有限个。直觉到此为止，剩下的由鸽巢原理接手。

\begin{center}
\begin{tikzpicture}[x=1cm,y=1cm,font=\small,>=stealth]
  \foreach \i in {0,1,2,3,4,5,6} {\draw[fill=blue!6] (\i,0) circle (0.18);}
  \foreach \i in {0,1,2,3,4,5} {\draw[->] (\i+0.22,0) -- (\i+0.78,0);}
  \draw[fill=orange!35] (2,0) circle (0.18);
  \draw[fill=orange!35] (4,0) circle (0.18);
  \node[font=\footnotesize] at (2,0) {\(s\)};
  \node[font=\footnotesize] at (4,0) {\(s\)};
  \draw[->] (4,0.22) to[bend right=45] (2,0.22);
  \node[font=\footnotesize] at (3,0.95) {绕这一圈，状态不变};
  \draw (-0.05,-0.35) -- (-0.05,-0.5) -- (1.95,-0.5) -- (1.95,-0.35);
  \node[below,font=\footnotesize] at (0.95,-0.5) {\(x\)};
  \draw (2.05,-0.35) -- (2.05,-0.5) -- (3.95,-0.5) -- (3.95,-0.35);
  \node[below,font=\footnotesize] at (3,-0.5) {\(y\)};
  \draw (4.05,-0.35) -- (4.05,-0.5) -- (6.05,-0.5) -- (6.05,-0.35);
  \node[below,font=\footnotesize] at (5.05,-0.5) {\(z\)};
\end{tikzpicture}

\emph{读一个长度至少为 \(\abs{Q}\) 的串，途中经过的状态比状态总数还多，必有一个状态被重复访问。两次访问之间那段输入 \(y\) 构成一个环：走 0 次、1 次或任意多次，终点都落在同一个状态上。机器分不清 \(xz\)、\(xyz\)、\(xyyz\)——它们对它来说是同一条经历。}
\end{center}

\begin{lemma}[泵引理]
若 \(L\) 正则，则存在正整数 \(p\)，使得每个 \(w\in L\) 且 \(\abs{w}\ge p\) 都能拆成 \(w=xyz\)，满足 \(\abs{xy}\le p\)、\(\abs{y}\ge1\)，且对一切 \(i\ge0\) 有 \(xy^iz\in L\)。
\end{lemma}

骨架就是上面那张图：取 \(p\) 为某台识别 \(L\) 的 DFA 的状态数，读 \(w\) 的前 \(p\) 个符号会经过 \(p+1\) 个状态，鸽巢原理保证其中两个相同；把第一次到达该状态之前的输入记为 \(x\)、两次之间的记为 \(y\)、其余记为 \(z\)，环可以绕任意多圈而终点不变，所以 \(xy^iz\) 与 \(w\) 同接受同拒绝。

用它否定 \(L=\set{a^nb^n}\)：假设 \(L\) 正则，得到抽水长度 \(p\)，取 \(w=a^pb^p\in L\)，长度 \(2p\ge p\)。由 \(\abs{xy}\le p\)，\(y\) 整段落在前 \(p\) 个符号里，因而全由 \(a\) 组成且非空。取 \(i=0\)，得到 \(xz=a^{p-\abs{y}}b^p\)，\(a\) 的个数严格少于 \(b\)，不在 \(L\) 中，与引理矛盾。故 \(L\) 非正则。

三个条件各挡一处漏洞。\(\abs{xy}\le p\) 把可泵段钉死在开头那 \(p\) 个符号里，正是这一条让我们断定 \(y\) 只含 \(a\)；没有它，对手可以把 \(y\) 放在 \(a\) 与 \(b\) 的交界上，泵出来的串未必违规。\(\abs{y}\ge1\) 排除空环，否则抽水什么也不改变，命题沦为废话。``对一切 \(i\ge0\)''把 \(i=0\) 也纳进来，上面的反证正是靠删除而不是靠重复完成的。

量词的顺序同样是条件的一部分：\(p\) 由对手给定，\(w\) 由你挑；分解 \(w=xyz\) 由对手给定，\(i\) 由你挑。所以挑 \(w\) 时不能挑短于 \(p\) 的串，讨论分解时也不能只讨论对自己方便的那一种。

\begin{remark}
泵引理只是必要条件。存在非正则语言照样满足抽水性质，因此``能泵''推不出正则，它只能用来否定。Myhill--Nerode 关系没有这个缺口：等价类个数有限当且仅当语言正则，两个方向都成立。用它处理 \(a^nb^n\) 甚至更快——前缀 \(a^0,a^1,a^2,\dots\) 两两可区分（\(a^i\) 与 \(a^j\) 用后缀 \(b^i\) 就分开了），等价类无限多，一句话结案。同一件事，泵引理讲的是``机器会混淆''，Myhill--Nerode 讲的是``语言本身要求的区分度''，后者更贴近问题。
\end{remark}

真要识别 \(a^nb^n\)，得给机器一块能随输入增长的存储。下推自动机加一个栈，读 \(a\) 压入一枚标记、读 \(b\) 弹出一枚，读完看栈空不空，这已经是上下文无关语言的地盘。回到开头那个正则表达式：括号配平与 \(a^nb^n\) 是同一类东西，任何嵌套深度无上界的结构都不在有限状态的射程内。该上解析器的时候就上解析器，再巧的写法也换不来一层额外记忆。

\subsection{把有限状态当成一份内存预算}

这条判断在流式处理里天天用得上。数据一遍流过、内存固定，能维护的统计量就受同一条限制约束：``目前为止 \(1\) 的个数是奇数还是偶数''只要两个状态；``当前是否处在一段引号内''也只要两个；而``目前为止 \(1\) 比 \(0\) 多''要求维护一个无上界的差值，有限状态办不到。想办到只有两条路，要么给出上界（差值截断在 \(\pm K\) 以内），要么承认误差换空间（近似计数）。设计前先问一句：哪些历史差异会改变未来的答案？答案有限，就能做成状态机。

同时要把语言类别与执行成本分开。一个语言是正则的，保证存在线性扫描的 DFA，不保证手头的实现高效：把 NFA 展开成 DFA 可能撞上指数状态数，而常见的回溯式正则引擎压根没有先构造 DFA，特定模式配上特定输入会跑出指数时间。``理论上可以线性''与``这个库现在很快''之间没有蕴含关系。

还有一个跨领域的呼应值得记一下，注意它是类比而不是定理。循环网络把历史压进一个固定维度的隐状态，这同样是一份预算；若把隐状态按有限精度量化，可区分的历史就只剩有限多种，输入足够长时两段不同的历史必然落进同一个向量，此后再也分不开——形状与上面那张鸽巢图完全一样。差别在于连续隐状态并非有限集合，训练出的表示也未必按``未来是否可区分''来分配容量，所以这里得不出泵引理式的硬结论。\hyperref[sec:14-Deep-Learning/05-Attention-and-Sequence-Models/01]{``Self-Attention 是内容决定的加权读取''}走的是另一条路：干脆不压缩，保留全部位置，按内容临时去取。定长摘要与全量保留之间的取舍，从有限自动机到序列模型一直在重复。

本单元两节共用一个念头：把可能无穷的东西压进有限的表示，再追问压缩之后还剩多少能力。上一节压的是同一时刻的条件组合，本节压的是历史。下一章\hyperref[ch:034]{``环与格：两种运算怎样组织代数与次序''}回到代数一侧，把``两种运算互相分配、互相吸收''这个模式单独抽出来；布尔代数会作为其中一个特例重新出现，那时它的运算律将从次序结构里长出来，不必再当作一张清单去记。
